\documentclass[11pt]{article}
\usepackage[a4paper,margin=27mm]{geometry}
\usepackage[T1]{fontenc}
\usepackage{lmodern,microtype,amsmath,amssymb,amsthm,mathtools,enumitem}
\usepackage[hidelinks,hypertexnames=false]{hyperref}
\newtheorem{theorem}{Theorem}
\newtheorem{lemma}[theorem]{Lemma}
\newtheorem{proposition}[theorem]{Proposition}
\newtheorem{corollary}[theorem]{Corollary}
\theoremstyle{remark}\newtheorem{remark}[theorem]{Remark}
\newcommand{\R}{\mathbb R}
\newcommand{\Dscr}{\mathcal D}
\newcommand{\Escr}{\mathcal E}
\newcommand{\supp}{\operatorname{supp}}
\newcommand{\one}{\mathbf 1}
\title{Finitely Additive Edge Dichotomy}
\author{}\date{10 October 2026}
\begin{document}\maketitle

\noindent
The finitely additive extension of the edge dichotomy suggested after
\emph{Edge Dichotomy and Leverage} (EDL) holds without an extra assumption
about splitting charges. It holds in every dimension, not just two.
The decisive observation is elementary: a finite nonnegative charge gives
arbitrarily small weight to \emph{some} positive-measure subset of any
positive-measure region. This is enough for continuity along an edge.

The resulting theorem says that a finitely additive primary is either
completely blind on mixed problems or has strictly positive \emph{ordinary}
increments there. Its purely finitely additive part need not vanish.
Consequently lexicographic hierarchies of finitely additive criteria cannot
give full coverage. We also extend EDL's two-additive-world theorem to
finitely additive one-sided extensions.

These results do not settle whether (G), (IB), and (S) exclude leverage
sequences for arbitrary rankings. None of the new arguments uses (IB).

\section{Setting}

There are $n\ge2$ worlds and $\Delta=\{p\ge0:\sum_i p_i=1\}$.
Binary problems are represented by $g\in H=\{g_1>0\}$, with normalized
Gaussian measure $\mu$. Coordinate hyperplanes are ignored. A model gives
one of four kinds of advice at each problem: recommend the better option,
remain silent, recommend the worse option, or reject both options. Their
ranks at world $i$ are $r_i=0,1,2,3$, respectively.

Models are measurable and identified modulo null sets. A model is
\emph{regular} if it never gives contradictory advice and is correct on
unanimous problems. Let $H'$ be the mixed problems. Every coherent model is
regular and has, up to a null set, the form
\[
 \Dscr_M=\{g:\underline{P_M}(g)>0\}\cup\{g\ge0:g\ne0\},
 \qquad \underline{P_M}(g)=\min_{m\in M}m\cdot g,
\]
for a nonempty compact convex $M\subseteq\Delta$. The ideal at $i$ is
$D_i=\{g:g_i>0\}$; the vacuous model is $V=\Dscr_\Delta$.

At each world $\preceq_i$ is a preorder; smaller means better.
\textbf{(G)} preserves pointwise improvement in guidance and makes it strict
when it occurs on a positive-measure set. \textbf{(S)} supplies coherent
undominated $C_p$, competitive on $\supp p$, with every section
\[
 A_i(X)=\{p:C_p\preceq_iX\}
\]
closed. Competitive means that against every measurable model $X$ there is
some supported $i$ with $C_p\preceq_iX$. Rankings need not be complete.

A \emph{regular primary} at $i$ is a real function $\Phi$ on regular models
such that $\Phi(X)<\Phi(Y)$ implies $X\prec_iY$. It follows that
$X\preceq_iY$ implies $\Phi(X)\le\Phi(Y)$ for regular $X,Y$.
We consider finitely additive primaries of the form
\begin{equation}
 \Phi(X)=a+\nu_1\{g\in H':r_i(g,X)\ge1\}
             +\nu_2\{g\in H':r_i(g,X)\ge2\},
 \label{eq:primary}
\end{equation}
where the $\nu_k$ are finite nonnegative finitely additive measures,
vanishing on $\mu$-null sets. They may vanish on positive-measure sets too.
The constant $a$ has no effect on comparisons.

We will use one consequence of coherence and competitiveness. On an edge
$[e_i,e_j]$, the credal set of $C_p$ lies on that edge. Otherwise an option
positive at both supported worlds, but sufficiently negative at an
unsupported world, would allow a strict guidance improvement at both
supported worlds. This contradicts competitiveness. In particular
$C_{e_i}=D_i$ and $C_{e_j}=D_j$.

Put $D_{ij}=\{g\in H':g_ig_j<0\}$ and $t=|g_i|/|g_j|$ on this set.
The edge credal set is an interval
$\{s e_i+(1-s)e_j:s\in[a_p,b_p]\}$. Set
$\alpha=(1-a_p)/a_p$ and $\beta=(1-b_p)/b_p$, allowing $0,\infty$.
Then $0\le\beta\le\alpha\le\infty$. At world $i$, $C_p$ is good for
$t>\alpha$, silent for $\beta<t<\alpha$, and bad for $t<\beta$.
Off $D_{ij}$ it is good at both $i$ and $j$.
Every finite level set of $t$ is null, and the distribution of $t$ on
$D_{ij}$ has a positive density on $(0,\infty)$.

\section{Finite additivity is enough for edge continuity}

\begin{lemma}[A cheap change always exists]\label{lem:cheap}
Let $\nu$ be a finite nonnegative finitely additive measure. If $\mu(Q)>0$
and $\epsilon>0$, there is $B\subseteq Q$ with $\mu(B)>0$ and
$\nu(B)<\epsilon$.
\end{lemma}
\begin{proof}
Partition $Q$ into $N$ positive-$\mu$-measure pieces, where
$N>\nu(Q)/\epsilon$. Finite additivity gives at least one piece with
charge at most $\nu(Q)/N<\epsilon$.
\end{proof}

This does not assert that every small-measure set has small charge, or
that a charge can be split into prescribed fractions. Neither assertion
is needed. We only need one cheap place to make a strict guidance change.

\begin{lemma}\label{lem:continuity}
Assume (G) and (S). If \eqref{eq:primary} is a regular primary at $i$, then
$f(p)=\Phi(C_p)$ is continuous on every edge $[e_i,e_j]$.
\end{lemma}
\begin{proof}
The ideal $D_i$ minimizes $\Phi$ on regular models. Among coherent edge
models, $D_j$ maximizes it, by the pointwise guidance order at $i$.

Suppose $p_m\to p$ on the edge and $f(p_m)\le f(p)-\epsilon$.
Then $C_p\ne D_i$. On a positive-measure subset of $D_{ij}$ its advice
can be strictly improved at $i$. Lemma~\ref{lem:cheap}, applied to
$\nu_1+\nu_2$, supplies such a subset on which the change has primary
cost less than $\epsilon$. The resulting regular $E$ satisfies
\[
 E\prec_i C_p,\qquad f(p)-\epsilon<\Phi(E)\le f(p).
\]
Thus $C_{p_m}\prec_i E$ for every $m$. Closedness of $A_i(E)$ gives
$C_p\preceq_i E$, a contradiction. The argument includes changes of
zero primary cost.

Conversely, suppose $f(p_m)\ge f(p)+\epsilon$. Then $C_p\ne D_j$.
Make a strict guidance improvement at $j$ on a positive-measure subset
of $D_{ij}$ whose $(\nu_1+\nu_2)$-charge is less than $\epsilon$.
The resulting regular $F$ satisfies
\[
 F\prec_j C_p,\qquad \Phi(F)<f(p)+\epsilon.
\]
Consequently $C_{p_m}\not\preceq_i F$. Competitiveness on the edge gives
$C_{p_m}\preceq_j F$, and closedness gives $C_p\preceq_j F$, again a
contradiction.
\end{proof}

\section{The strict-or-blind alternative}

Decompose the increments into ordinary and purely finitely additive parts:
\begin{equation}
 \nu_1=c\mu+\pi_1,\qquad \nu_2=d\mu+\pi_2,\qquad c,d\ge0.
 \label{eq:split}
\end{equation}
Here $c,d$ are integrable. The Yosida--Hewitt decomposition supplies such a
decomposition and concentrating sets $Z_m$ with $\mu(Z_m)\to0$ and
$\pi_k(H'\setminus Z_m)=0$ \cite[Theorem 1 and its proof]{PP}.
In particular, if $c=0$ on a positive-measure set $Q$, then some
positive-measure $B\subseteq Q$ has $\nu_1(B)=0$: take
$B=Q\setminus Z_m$ for sufficiently large $m$. The corresponding statement
holds for $d$ and $\nu_2$.

\begin{theorem}[Finitely additive edge dichotomy]\label{thm:dichotomy}
Assume (G) and (S), and let \eqref{eq:primary} be a regular primary at $i$.
For every $j\ne i$, exactly one alternative holds:
\begin{enumerate}[label=(\alph*),leftmargin=*]
\item $c>0$ and $d>0$ almost everywhere on $D_{ij}$;
\item $\nu_1(D_{ij})=\nu_2(D_{ij})=0$.
\end{enumerate}
If $n\ge3$, the alternatives agree for all $j\ne i$. Thus either the
ordinary parts of both increments are strictly positive a.e. on $H'$, or
the entire primary is constant on regular models.
\end{theorem}
\begin{proof}
Fix $j$. Write $m_0=\Phi(D_i)$ and
\[
 J(x)=(\nu_1+\nu_2)(D_{ij}\cap\{t<x\}),\qquad 0\le x\le\infty.
\]
For a selected edge model with thresholds $\alpha\ge\beta$,
\begin{equation}
 f(p)=m_0+\nu_1(D_{ij}\cap\{t<\alpha\})
             +\nu_2(D_{ij}\cap\{t<\beta\}),\qquad
 m_0+J(\beta)\le f(p)\le m_0+J(\alpha).
 \label{eq:bounds}
\end{equation}
By Lemma~\ref{lem:continuity}, $f$ takes every value in
$[m_0,m_0+J(\infty)]$ on the edge. No continuity of $J$ is assumed.

\emph{Invisible changes force local extrema.}
If $E\prec_i C_p$ is regular and $\Phi(E)=f(p)$, then closedness excludes
$C_q\preceq_i E$ near $p$. The primary property gives $f(q)\ge f(p)$:
$p$ is a local minimum. If instead $F\prec_j C_p$ and $\Phi(F)=f(p)$,
closedness excludes $C_q\preceq_j F$ near $p$ on the edge.
Competitiveness gives $C_q\preceq_i F$, so $f(q)\le f(p)$:
$p$ is a local maximum on the edge.

Apply the zero-charge observation following \eqref{eq:split} to make these
changes wherever the relevant density vanishes. This gives four sufficient
conditions:
\begin{center}
\begin{tabular}{ccl}
Density zero on a non-null set & Location & Consequence\\[2pt]
$c=0$ & $\beta<t<\alpha$ & local minimum: silence to $i$-good\\
$d=0$ & $t<\beta$ & local minimum: $i$-bad to silence\\
$c=0$ & $t>\alpha$ & local maximum: $j$-bad to silence\\
$d=0$ & $\beta<t<\alpha$ & local maximum: silence to $j$-good
\end{tabular}
\end{center}
All sets in this table are restricted to $D_{ij}$.
The values at local minima of a real function on a second-countable space
form a countable set: assign to each value a basic neighbourhood on which
it is the infimum. Distinct values require distinct neighbourhoods.
The same argument applies to local maxima. Thus the values accounted for
by the table form a countable set.

\emph{Every other value lies between two fixed thresholds.}
Call a selected point free if none of the four conditions applies. Put
\[
 U=\inf\{x:\mu(D_{ij}\cap\{c=0\}\cap\{t>x\})=0\},
 \qquad
 L=\sup\{x:\mu(D_{ij}\cap\{d=0\}\cap\{t<x\})=0\},
\]
where $x\in[0,\infty]$. Continuity of the ordinary measure $\mu$ gives
\begin{equation}
 c>0\ \text{a.e. on }D_{ij}\cap\{t>U\},\qquad
 d>0\ \text{a.e. on }D_{ij}\cap\{t<L\}.
 \label{eq:positive-tails}
\end{equation}
At a free point, the table and the nullity of threshold level sets give
$U\le\beta\le\alpha\le L$. Its value therefore lies in
$[m_0+J(U),m_0+J(L)]$.

If there are no free points, the entire interval of values is countable,
so $J(\infty)=0$, giving (b). Otherwise $U\le L$, and every value in
\[
 [m_0,m_0+J(U))\ \cup\ (m_0+J(L),m_0+J(\infty)]
\]
is a local-extremum value. Both intervals must be empty. Hence
\begin{equation}
 J(U)=0,\qquad J(L)=J(\infty).
 \label{eq:tails-zero}
\end{equation}
Both charges, and therefore both densities, vanish on the tails $t<U$
and $t>L$. This uses only finite additivity and the fact that $\{t=L\}$
is null when $L$ is finite.

If $U>0$ and $L>0$, the region $0<t<\min(U,L)$ has positive measure.
There $d=0$ by \eqref{eq:tails-zero}, but $d>0$ by
\eqref{eq:positive-tails}, a contradiction. If $L=0$, then
$J(\infty)=J(0)=0$, giving (b). In the remaining case $U=0$,
\eqref{eq:positive-tails} gives $c>0$ a.e. throughout $D_{ij}$.
Then \eqref{eq:tails-zero} forces $L=\infty$, and
\eqref{eq:positive-tails} gives $d>0$ a.e. too. This is (a).

For $n\ge3$, the regions $D_{ij}$ cover $H'$ and any two with fixed $i$
overlap on a positive-measure set. The strict and blind alternatives cannot
hold on overlapping regions. If they are all blind, finite subadditivity
gives $\nu_1(H')=\nu_2(H')=0$.
\end{proof}

\begin{remark}
The strict alternative does not remove the purely finitely additive parts.
CBOI's Proposition 5 supplies a primary with strictly positive ordinary
increments and nonzero pure increments, satisfying (S) but having no ordinary
additive weak extension. Accordingly, $\Phi$ itself supplies a finitely
additive weak extension; its ordinary part need not supply a weak extension.
\end{remark}

\section{Lexicographic hierarchies of any length}

The obstruction from CBOI, strengthened by EDL's Proposition 1, says:
if each world's comparisons from coherent undominated models to regular
models have a strictly guidance-respecting finitely additive weak extension,
then, for $n\ge3$, every target containing three simplex vertices is
dominated by a coherent model. We apply that result here.

There is also a local conclusion. At any world with a leverage sequence,
every finitely additive regular primary must be constant on regular models.
Indeed, a nonconstant primary is strict by Theorem~\ref{thm:dichotomy}
(also when $n=2$, where $H'=D_{12}$). Its increments give a strictly
positive finitely additive extension, contradicting EDL's exact leverage
test. Thus such a world cannot even retain a nontrivial finitely additive
criterion as an overriding priority.

\begin{corollary}\label{cor:lex}
Let $n\ge3$, and assume (G), (LA), and (S). At each world suppose the ranking
on regular models respects a lexicographic hierarchy of criteria of the
form \eqref{eq:primary}, and at least one criterion is nonconstant there.
The hierarchy may be of any length, provided every unequal pair of criterion
vectors has a first differing index. Then every $\Dscr_M$ whose credal set
contains three simplex vertices is dominated by a coherent model.
\end{corollary}
\begin{proof}
At world $i$, let $W_i$ be the regular model that gives bad advice at $i$
on every mixed problem. Compare its criterion vector with that of $D_i$.
Exactly the nonconstant criteria distinguish this pair, because their
increments are nonnegative. Thus the first differing index is the first
nonconstant criterion, even if the index order is not well ordered.

All preceding criteria are constant on regular models. The first
nonconstant criterion is therefore a regular primary. By
Theorem~\ref{thm:dichotomy}, both its ordinary increments are positive
a.e. on $H'$. Its finitely additive increments are consequently strictly
positive on every positive-measure subset of $H'$.
Extending the increments positively on unanimous problems, and adding a
positive contradiction increment, gives a strictly guidance-respecting
finitely additive weak extension on regular models. Apply the stated
obstruction at all worlds.
\end{proof}

This includes finite, countable, and transfinite well-ordered hierarchies.
If the hierarchy itself determines all comparisons, strict (G) ensures
that it is nonconstant at every world. If arbitrary extra tie-breaking is
allowed, a world at which all criteria are constant merely delegates the
entire problem to that extra rule.

\section{Two finitely additive worlds and one unrestricted world}

EDL's Theorem 8 also extends, but its conclusion must be expressed using
the ordinary parts of the two finitely additive losses. We do not assert
that those ordinary parts weakly extend the original comparisons.

\begin{proposition}\label{prop:two-worlds}
Let $n=3$, assume (G), (LA), and (S), and suppose $V$ is undominated.
At worlds $2,3$, suppose the comparisons from coherent undominated models
to regular models have strictly guidance-respecting finitely additive weak
extensions $L_2,L_3$. Write $L_k=L_k^c+L_k^s$ for their ordinary and pure
parts, and let $c_k,b_k$ be the silence and bad penalties of $L_k^c$.
There is $\kappa^*\in(0,\infty)$ such that, on problems oriented with
$g_2>0>g_3$,
\[
 \frac{c_2}{b_2}+\frac{c_3}{b_3}\le1,\qquad
 c_2\le\kappa^*(b_3-c_3),\qquad
 \kappa^*c_3\le b_2-c_2
 \quad\text{a.e.}
\]
Thus silence minimizes ordinary expected loss at weights $(1,\kappa^*)$
on every $(2,3)$-mixed problem. World $1$ needs no numerical representation.
\end{proposition}
\begin{proof}
Strict positivity of each finitely additive increment implies that its
ordinary density is positive a.e. Indeed, on a positive-measure set where
the density vanished, removing a sufficiently small concentrating set for
the pure part would leave a positive-measure set of zero total charge.
Thus $L_2^c,L_3^c$ have strictly ordered local penalties.

The proof of EDL's Theorem 8 needs two ingredients. We establish both for
$L^c$, using common concentrating sets $Z_m$ for the pure parts, with
$\mu(Z_m)\to0$.

\emph{Coherent ordinary Bayes models on the edge.} Fix $\kappa>0$.
Let $B$ minimize $L_2^c+\kappa L_3^c$ pointwise, with minimum $a$.
Let regular $E_m$ approach the infimum $b$ of $L_2^s+\kappa L_3^s$.
Splice $E_m$ on $Z_m$ with $B$ elsewhere to obtain regular $X_m$.
Then $(L_2+\kappa L_3)(X_m)\to a+b$.

The closed sets $A_2(X_m)$ and $A_3(X_m)$ cover the connected edge
$[e_2,e_3]$ and contain its respective endpoints. They intersect, giving
a selected $C_m$ weakly better than $X_m$ at both worlds. The one-sided
extension inequalities imply
\[
 a\le(L_2^c+\kappa L_3^c)(C_m)
 \le(L_2+\kappa L_3)(X_m)-b\longrightarrow a.
\]
The edge credal intervals have a Hausdorff-convergent subsequence; their
advice converges a.e. to the limiting coherent edge model, since boundary
hyperplanes are null. Ordinary integral losses converge too. The limit
is a coherent Bayes model for $L_2^c+\kappa L_3^c$.

\emph{No trade improving world 1 and both ordinary losses.}
Let regular $X$ agree with $V$ except for recommending an option positive
at world $1$ on a positive-measure set of mixed problems. Suppose both
$L_2^c(X)<L_2^c(V)$ and $L_3^c(X)<L_3^c(V)$.
Replace $X$ by $V$ on $Z_m$, obtaining $Y_m$. For large $m$, a
positive-measure part of the recommendation remains, so $Y_m\prec_1V$
by (G). Meanwhile
\[
 L_k(Y_m)-L_k(V)
   =L_k^c(Y_m)-L_k^c(V)
   \longrightarrow L_k^c(X)-L_k^c(V)<0,\qquad k=2,3.
\]
The undominated $V$ is competitive on some support by the simplex lemma,
and hence on all three worlds. It cannot satisfy $V\preceq_1Y_m$; the
one-sided extension inequalities rule out $V\preceq_kY_m$ for $k=2,3$.
This is a contradiction.

These are precisely Steps 1 and 3 needed in EDL's proof, with $L^c$ in
place of $L$. Steps 2 and 4--6 then use only the ordinary penalties and
atomless splitting: the two Bayes thresholds are monotone functions of
$g_2/(-g_3)$, and the no-trade condition separates the cross-world ratios.
They give the three displayed inequalities. No extension property of
$L^c$ has been used.
\end{proof}

Consequently EDL's Corollary 9 applies to the contrasts of these ordinary
parts. For every finite sign-preserving world-$1$ contrast $\psi_1$, some
prior with $p_1>0$ admits no coherent contrast completion. This follows by
the same three-world compatibility argument as in that corollary.

\section{What is still missing}

EDL's leverage test is an exact test for a one-sided finitely additive
extension. Its use of Rokhlin's norm-topology Kreps--Yan theorem is valid.
The weak-star variant supplies the corresponding test for ordinary
extensions. These tests should not be confused with the stronger question
of constructing a ranking satisfying all the axioms and leaving $V$
undominated: a ranking with leverage alone need not do that.

The present results settle the proposed finitely additive edge sub-question
and rule out lexicographic finitely additive hierarchies of arbitrary length.
They also extend the two-additive-world reduction. They do not show that
an arbitrary independent ranking admits any such hierarchy or any additive
primary. That representation step remains unproved. In particular, the
argument does not exclude leverage sequences for genuinely non-additive
incomplete rankings.

\begin{thebibliography}{9}
\bibitem{PP} T.~Pennanen and A.-P.~Perkki\"o,
Topological duals of locally convex function spaces,
\emph{Positivity} 26, article 2 (2022), Theorem 1.
\href{https://doi.org/10.1007/s11117-022-00867-4}{doi:10.1007/s11117-022-00867-4}.
\end{thebibliography}
\end{document}
